\section{DESTINO Alignment and Evidence Boundary}
\label{sec:destino_context}

The original DESTINO proposal is the reference deployment context for this thesis. It describes a broad decentralized
ecosystem in which software-quality and maintenance evidence contributes to developer reputation through automated
measurement, signed submissions, immutable ledger records, evidence storage and discovery, and developer-oriented
profiles~\cite{destinoProposal}.

This thesis intentionally implements a narrower part of that vision. It compares repository-derived candidate indicators,
selects one signal for a baseline proof of concept, and evaluates whether its evidence can be calculated deterministically,
submitted by an authorized party, committed with integrity protections, and independently replayed. It does not implement
qualified electronic signatures, maintenance-response evidence, aggregated reputation profiles, public search services, or
production governance. Those broader DESTINO elements remain outside the claims evaluated here.

DESTINO is therefore not treated as a reason to move all computation on-chain. Instead, it defines an evidence boundary that
every candidate metric must satisfy before a proof-of-concept metric can be selected. Within that boundary, authorization
controls who may submit metric evidence; the selected metric does not itself grant or deny developer permissions.

The working DESTINO-aligned evidence model has four parts:
\begin{itemize}[leftmargin=*]
  \item a developer-facing submission containing project/archive metadata, software-type tags, storage-persistence metadata,
  and a cryptographic commitment to the submitted artifact;
  \item an off-chain measurement pipeline that retrieves the committed artifact, verifies integrity, and computes the selected metric;
  \item compact on-chain records containing the submission reference, metric value, and evidence hash;
  \item auditor-accessible replay: retrieve the same artifact, rerun the declared script and configuration, and compare the result
  with the committed metric record.
\end{itemize}

This boundary keeps expensive repository traversal, Git history analysis, and tabular validation off-chain, while the smart-contract
layer stores only bounded metadata, commitments, and metric outputs. This matches common blockchain engineering constraints:
large artifacts and unbounded loops are poor on-chain candidates, while hash commitments and compact records are auditable and
storage-bounded~\cite{blockchainsurvey,smartmetrics2023}.

\subsection{Implications for This Thesis}
\label{sec:destino_phase1_implications}

Phase I is therefore not solely a data-extraction phase. It prepares the evidence artifacts subsequently used by
the DESTINO-aligned PoC. The alignment is summarized in Table~\ref{tab:destino_alignment}.

\begin{table}[H]
\centering
\footnotesize
\renewcommand{\arraystretch}{1.15}
\begin{tabularx}{\textwidth}{p{3.4cm}X X}
\toprule
\textbf{DESTINO need} & \textbf{Phase I artifact or rule} & \textbf{Why it matters} \\
\midrule
Stable artifact identity &
\path{data/dataset_freeze_all.csv}, \path{data/versions_all.csv}, frozen commit hashes &
Metric evidence can be tied to explicit projects, tags, and repository states. \\
Replayable measurement interval &
\path{data/windows_all.csv} &
Each metric value has a deterministic release-window boundary. \\
Compact metric evidence &
\path{out/loc_per_tag_all.csv}, clean metric tables, normalized tables, human-only M3/M4 tables &
The contract can store selected values without storing raw repositories or diffs. \\
Audit and validation evidence &
\path{logs/*.log}, cleaning checks, duration flags, bot-sensitivity summaries &
A verifier can inspect how raw extraction became analysis-ready evidence. \\
Smart-contract selection discipline &
Criteria C1--C4 plus the feasibility rubric &
The selected PoC metric must be empirically useful and feasible under bounded contract execution. \\
\bottomrule
\end{tabularx}
\caption{Alignment between DESTINO evidence requirements and current Phase I artifacts.}
\label{tab:destino_alignment}
\end{table}

Two Phase I design decisions are especially important for DESTINO alignment. First, duration-normalized outputs and
\texttt{sensitivity\_only} flags prevent release cadence artifacts from being silently recorded as comparable metric evidence.
Second, human-only M3/M4 outputs prevent automated accounts from being interpreted as human maintainer behavior. The all-account
tables remain useful as an audit trail, but the primary developer-behavior evidence is explicitly human-only.

No PoC metric is selected within Phase I itself. At that phase boundary, DESTINO alignment means the pipeline produces evidence
that can subsequently be committed, audited, and scored; Phase II performs the comparative and feasibility gate.
